\documentclass[12pt,a4paper]{article}

\usepackage[utf8]{inputenc}
\usepackage[T2A]{fontenc}
\usepackage[russian]{babel}
\usepackage{amsmath,amssymb}
\usepackage{geometry}
\usepackage{graphicx}
\usepackage{booktabs}
\usepackage{float}
\usepackage{array}
\usepackage{caption}

\geometry{margin=2.2cm}

\newcommand{\plotplaceholder}[2][0.92\textwidth]{%
    \IfFileExists{#2}{%
        \includegraphics[width=#1]{#2}%
    }{%
        \fbox{%
            \begin{minipage}[c][0.24\textheight][c]{#1}
            \centering
            Заглушка для графика\\[4pt]
            \texttt{\detokenize{#2}}
            \end{minipage}%
        }%
    }%
}

\title{Отчёт по задачам 2, 3 и 4\\[4pt]
Оценивание плотности распределения Вейбулла--Гнеденко}
\author{}
\date{}

\begin{document}

\maketitle

\section*{1. Условие задачи}

Рассматривается генеральная совокупность с плотностью распределения Вейбулла--Гнеденко
\[
f(x)=\lambda c x^{c-1} e^{-\lambda x^c}, \qquad \lambda>0,\; c>0,\; x\ge 0.
\]

Во всех трёх задачах используются параметры
\[
\lambda=1, \qquad c=3, \qquad n=400.
\]

Следовательно,
\[
f(x)=3x^2e^{-x^3}, \qquad x\ge 0.
\]

\textbf{Задача 2.} Построить гистограммную оценку плотности и график зависимости средней интегральной квадратической ошибки $\overline{\delta_n}(m)$ от числа разрядов $m$ для значений
\[
m=5(5)25 \quad \text{и} \quad m=10(10)40.
\]

\textbf{Задача 3.} Построить проекционную оценку плотности по функциям Лагерра и исследовать зависимость $\overline{\delta_n}(N)$ от числа членов разложения $N$.

\textbf{Задача 4.} Построить ядерную оценку плотности с ядром Епанечникова
\[
\widetilde{f}_h(x)=\frac{1}{nh}\sum_{i=1}^{n}K\!\left(\frac{x-x_i}{h}\right),
\]
где
\[
K(u)=
\begin{cases}
\dfrac{3}{4\sqrt{5}}\left(1-\dfrac{u^2}{5}\right), & |u|\le\sqrt{5},\\[6pt]
0, & |u|>\sqrt{5}.
\end{cases}
\]
Требуется вычислить ОИСКО и исследовать зависимость ошибки от ширины окна $h$.

\section*{2. Кратко теория}

Выборка строится методом Монте-Карло через метод обращения. Для распределения Вейбулла--Гнеденко при $\lambda=1$ и $c=3$
\[
F(x)=1-e^{-x^3}.
\]
Если $U\sim U(0,1)$, то
\[
X=F^{-1}(U)=\sqrt[3]{-\ln(1-U)}.
\]

Во всех задачах качество оценки плотности измеряется интегральной квадратической ошибкой
\[
\delta_n=\int_0^{\infty}\bigl(\widetilde{f}(x)-f(x)\bigr)^2dx.
\]
Так как выборка случайна, ошибка тоже случайна, поэтому оценивается её среднее по многим прогонам:
\[
\overline{\delta_n}=\frac{1}{R}\sum_{r=1}^{R}\delta_n^{(r)}.
\]

В задаче 2 используется гистограммная оценка,
в задаче 3 -- проекционная оценка по функциям Лагерра,
в задаче 4 -- ядерная оценка с ядром Епанечникова.

\section*{3. Идея решения и алгоритм}

Для всех трёх задач используется один и тот же сценарий:
\begin{enumerate}
    \item сгенерировать выборку $x_1,\dots,x_n$ методом обращения;
    \item построить нужный тип оценки плотности;
    \item вычислить интегральную квадратическую ошибку;
    \item повторить моделирование много раз;
    \item усреднить значения ошибки и построить график зависимости от параметра модели.
\end{enumerate}

Все итоговые графики сохраняются кодом в папку \texttt{plots/}. В этом варианте оформления:
\begin{itemize}
    \item в задаче 2 гистограммы сгруппированы в две сводные фигуры: для набора $m=5(5)25$ используется раскладка 2 графика в ширину и 3 в высоту, а для набора $m=10(10)40$ -- раскладка $2\times 2$;
    \item в задаче 3 выбранные проекционные оценки для разных $N$ собраны в одну сводную фигуру с раскладкой 2 графика в ширину и 3 в высоту;
    \item в задаче 4 четыре ядерные оценки для разных $h$ собраны в одну фигуру $2\times 2$, а оценка при лучшем $h$ вынесена отдельно.
\end{itemize}

\section*{4. Результаты}

\subsection*{4.1. Задача 2. Гистограмма}

\begin{table}[H]
\centering
\caption{Средние значения $\overline{\delta_n}(m)$}
\begin{tabular}{ccc}
\toprule
$m$ & $\overline{\delta_n}(m)$ & стандартное отклонение \\
\midrule
5  & 0.084073 & 0.002462 \\
10 & 0.033474 & 0.004856 \\
15 & 0.022548 & 0.006532 \\
20 & 0.021594 & 0.007747 \\
25 & 0.023139 & 0.008878 \\
30 & 0.025315 & 0.009652 \\
40 & 0.033219 & 0.011598 \\
\bottomrule
\end{tabular}
\end{table}

Минимум ошибки достигается примерно при
\[
m \approx 20.
\]

\begin{figure}[H]
\centering
\plotplaceholder{plots/task2_hist_grid_5_25_step5.png}
\caption{Гистограммные оценки для набора $m=5(5)25$}
\end{figure}

\begin{figure}[H]
\centering
\plotplaceholder{plots/task2_hist_grid_10_40_step10.png}
\caption{Гистограммные оценки для набора $m=10(10)40$}
\end{figure}

\begin{figure}[H]
\centering
\plotplaceholder{plots/task2_delta_vs_m.png}
\caption{Зависимость $\overline{\delta_n}(m)$ от числа разрядов $m$}
\end{figure}

\subsection*{4.2. Задача 3. Проекционные оценки}

\begin{table}[H]
\centering
\caption{Средние значения $\overline{\delta_n}(N)$}
\begin{tabular}{ccc}
\toprule
$N$ & $\overline{\delta_n}(N)$ & стандартное отклонение \\
\midrule
10  & 0.095576 & 0.001290 \\
20  & 0.019040 & 0.001875 \\
30  & 0.007125 & 0.002438 \\
40  & 0.006095 & 0.003161 \\
50  & 0.004759 & 0.003381 \\
60  & 0.005017 & 0.003870 \\
70  & 0.005562 & 0.004229 \\
80  & 0.005802 & 0.004208 \\
90  & 0.006632 & 0.004910 \\
100 & 0.006850 & 0.004618 \\
\bottomrule
\end{tabular}
\end{table}

Минимум ошибки достигается примерно при
\[
N \approx 50.
\]

\begin{figure}[H]
\centering
\plotplaceholder{plots/task3_proj_grid.png}
\caption{Проекционные оценки для выбранных значений $N$}
\end{figure}

\begin{figure}[H]
\centering
\plotplaceholder{plots/task3_delta_vs_N.png}
\caption{Зависимость $\overline{\delta_n}(N)$ от числа членов разложения $N$}
\end{figure}

\subsection*{4.3. Задача 4. Ядерные оценки}

\begin{table}[H]
\centering
\caption{Средние значения $\overline{\delta_n}(h)$}
\begin{tabular}{ccc}
\toprule
$h$ & $\overline{\delta_n}(h)$ & стандартное отклонение \\
\midrule
0.05 & 0.011493 & 0.005370 \\
0.10 & 0.005725 & 0.003682 \\
0.15 & 0.007410 & 0.004540 \\
0.20 & 0.015862 & 0.005199 \\
0.25 & 0.030716 & 0.005361 \\
0.30 & 0.052859 & 0.005411 \\
0.35 & 0.080748 & 0.004876 \\
0.40 & 0.113674 & 0.004004 \\
0.45 & 0.147071 & 0.003142 \\
0.50 & 0.181734 & 0.002672 \\
0.55 & 0.214935 & 0.002129 \\
0.60 & 0.246198 & 0.001997 \\
0.65 & 0.275902 & 0.001960 \\
0.70 & 0.302357 & 0.001877 \\
0.75 & 0.327212 & 0.001641 \\
0.80 & 0.349962 & 0.001710 \\
0.85 & 0.370880 & 0.001536 \\
0.90 & 0.390419 & 0.001489 \\
0.95 & 0.408411 & 0.001364 \\
1.00 & 0.425324 & 0.001392 \\
\bottomrule
\end{tabular}
\end{table}

На исследованной сетке минимум достигается примерно при
\[
h \approx 0.10.
\]

Для визуального сравнения дополнительно рассмотрены четыре значения $h$: $0.05$, $0.10$, $0.20$ и $0.50$.

\begin{figure}[H]
\centering
\plotplaceholder{plots/task4_delta_vs_h.png}
\caption{Зависимость $\overline{\delta_n}(h)$ от ширины окна $h$}
\end{figure}

\begin{figure}[H]
\centering
\plotplaceholder{plots/task4_kde_grid.png}
\caption{Ядерные оценки для нескольких значений ширины окна $h$}
\end{figure}

\begin{figure}[H]
\centering
\plotplaceholder{plots/task4_kde_best_h.png}
\caption{Ядерная оценка при значении $h$, близком к оптимальному}
\end{figure}

\subsection*{4.4. Общий вывод}

Полученные результаты показывают, что для каждого способа восстановления плотности существует свой параметр, при котором средняя интегральная квадратическая ошибка минимальна:
\[
m \approx 20, \qquad N \approx 50, \qquad h \approx 0.10.
\]

\end{document}
